% GENERATED FILE. Edit canonical lessons, then run npm run generate:books.
% canonical-source-hash: fnv1a64:b94a696c925775c7
% canonical-lessons: AR-C221-taqrir, AR-C221-ijtima, AR-C221-aqd, AR-C221-imtihan, AR-C221-wajib

\chapter{Report, Meeting, Contract, Exam, Homework}
\label{ch:ar-report-meeting-contract-exam-homework}
\hlchaptermodality{221}

\begin{chapteropening}
\textbf{By the end of this chapter:} \emph{I can say a report, a meeting, a contract, an exam and homework.}

Complete the last lesson of chapter 221: I can say a report, a meeting, a contract, an exam and homework.
\end{chapteropening}

\section[\texorpdfstring{\ar{تقرير} (taqrīr)}{تقرير (taqrīr)}]{\ar{تقرير} (taqrīr) — a report}

\label{lesson:AR-C221-taqrir}



\begin{quote}
Before the new one: say the Arabic for an account, then the Arabic for a butcher.
\end{quote}

\subsection*{You'll want to know: \ar{تقرير}}
\textbf{\ar{تقرير}} — \emph{taqrīr} — \textquotedblleft{}a report\textquotedblright{}.

\begin{cousinweb}
\textbf{\ar{تقرير}} (\emph{taqrīr}) is built on the native Arabic root \emph{q-r-r}, \textquotedblleft{}to settle, to be still\textquotedblright{} (Lane). \emph{Taqrīr} is the verbal noun of Form II, \textquotedblleft{}establishing\textquotedblright{}, hence \textquotedblleft{}a report\textquotedblright{}.
\end{cousinweb}

\begin{grammarlens}[title={What you've built}]
One more word for this chapter.
\end{grammarlens}

\subsection*{Your turn}
\begin{itemize}
\raggedright
  \item \emph{Say it:} \emph{taqrīr}
  \item \emph{Say it:} \emph{taqrīr}, once more
  \item \emph{Say it:} say \emph{taqrīr}
  \item \emph{Recall:} say the Arabic for an account, then the Arabic for a butcher, then say \emph{taqrīr} again
\end{itemize}

\begin{tcolorbox}[breakable,colback=teal!4,colframe=teal!35!black,title={Before you move on}]
What does \ar{تقرير} mean? (\textquotedblleft{}A report\textquotedblright{}.) Say it once more.
\end{tcolorbox}
\section[\texorpdfstring{\ar{اجتماع} (ijtimāʿ)}{اجتماع (ijtimāʿ)}]{\ar{اجتماع} (ijtimāʿ) — a meeting}

\label{lesson:AR-C221-ijtima}



\begin{quote}
Before the new one: say the Arabic for a butcher, then the Arabic for a report.
\end{quote}

\subsection*{You'll want to know: \ar{اجتماع}}
\textbf{\ar{اجتماع}} — \emph{ijtimāʿ} — \textquotedblleft{}a meeting\textquotedblright{}.

\begin{cousinweb}
\textbf{\ar{اجتماع}} (\emph{ijtimāʿ}) is built on the native Arabic root \emph{j-m-ʿ}, \textquotedblleft{}to gather\textquotedblright{} (Lane). \emph{Ijtimāʿ} is the verbal noun of Form VIII, \textquotedblleft{}coming together\textquotedblright{}.
\end{cousinweb}

\begin{grammarlens}[title={What you've built}]
One more word for this chapter.
\end{grammarlens}

\subsection*{Your turn}
\begin{itemize}
\raggedright
  \item \emph{Say it:} \emph{ijtimāʿ}
  \item \emph{Say it:} \emph{ijtimāʿ}, once more
  \item \emph{Say it:} say \emph{ijtimāʿ}
  \item \emph{Recall:} say the Arabic for a butcher, then the Arabic for a report, then say \emph{ijtimāʿ} again
\end{itemize}

\begin{tcolorbox}[breakable,colback=teal!4,colframe=teal!35!black,title={Before you move on}]
What does \ar{اجتماع} mean? (\textquotedblleft{}A meeting\textquotedblright{}.) Say it once more.
\end{tcolorbox}
\section[\texorpdfstring{\ar{عقد} (ʿaqd)}{عقد (ʿaqd)}]{\ar{عقد} (ʿaqd) — a contract}

\label{lesson:AR-C221-aqd}



\begin{quote}
Before the new one: say the Arabic for a report, then the Arabic for a meeting.
\end{quote}

\subsection*{You'll want to know: \ar{عقد}}
\textbf{\ar{عقد}} — \emph{ʿaqd} — \textquotedblleft{}a contract\textquotedblright{}.

\begin{cousinweb}
\textbf{\ar{عقد}} (\emph{ʿaqd}) is built on the native Arabic root \emph{ʿ-q-d}, \textquotedblleft{}to tie a knot\textquotedblright{} (Lane). \emph{ʿAqd} is a tie, a binding: a contract.
\end{cousinweb}

\begin{grammarlens}[title={What you've built}]
One more word for this chapter.
\end{grammarlens}

\subsection*{Your turn}
\begin{itemize}
\raggedright
  \item \emph{Say it:} \emph{ʿaqd}
  \item \emph{Say it:} \emph{ʿaqd}, once more
  \item \emph{Say it:} say \emph{ʿaqd}
  \item \emph{Recall:} say the Arabic for a report, then the Arabic for a meeting, then say \emph{ʿaqd} again
\end{itemize}

\begin{tcolorbox}[breakable,colback=teal!4,colframe=teal!35!black,title={Before you move on}]
What does \ar{عقد} mean? (\textquotedblleft{}A contract\textquotedblright{}.) Say it once more.
\end{tcolorbox}
\section[\texorpdfstring{\ar{امتحان} (imtiḥān)}{امتحان (imtiḥān)}]{\ar{امتحان} (imtiḥān) — an exam}

\label{lesson:AR-C221-imtihan}



\begin{quote}
Before the new one: say the Arabic for a meeting, then the Arabic for a contract.
\end{quote}

\subsection*{You'll want to know: \ar{امتحان}}
\textbf{\ar{امتحان}} — \emph{imtiḥān} — \textquotedblleft{}an exam\textquotedblright{}.

\begin{cousinweb}
\textbf{\ar{امتحان}} (\emph{imtiḥān}) is built on the native Arabic root \emph{m-ḥ-n}, \textquotedblleft{}to try, to test\textquotedblright{} (Lane). \emph{Imtiḥān} is the verbal noun of Form VIII, \textquotedblleft{}testing\textquotedblright{}.
\end{cousinweb}

\begin{grammarlens}[title={What you've built}]
One more word for this chapter.
\end{grammarlens}

\subsection*{Your turn}
\begin{itemize}
\raggedright
  \item \emph{Say it:} \emph{imtiḥān}
  \item \emph{Say it:} \emph{imtiḥān}, once more
  \item \emph{Say it:} say \emph{imtiḥān}
  \item \emph{Recall:} say the Arabic for a meeting, then the Arabic for a contract, then say \emph{imtiḥān} again
\end{itemize}

\begin{tcolorbox}[breakable,colback=teal!4,colframe=teal!35!black,title={Before you move on}]
What does \ar{امتحان} mean? (\textquotedblleft{}An exam\textquotedblright{}.) Say it once more.
\end{tcolorbox}
\section[\texorpdfstring{\ar{واجب} (wājib)}{واجب (wājib)}]{\ar{واجب} (wājib) — homework; a duty}

\label{lesson:AR-C221-wajib}



\begin{quote}
Before the new one: say the Arabic for a contract, then the Arabic for an exam.
\end{quote}

\subsection*{You'll want to know: \ar{واجب}}
\textbf{\ar{واجب}} — \emph{wājib} — \textquotedblleft{}homework; a duty\textquotedblright{}.

\begin{cousinweb}
\textbf{\ar{واجب}} (\emph{wājib}) is built on the native Arabic root \emph{w-j-b}, \textquotedblleft{}to be due, to fall\textquotedblright{} (Lane). \emph{Wājib} is the active participle, \textquotedblleft{}due, required\textquotedblright{}.
\end{cousinweb}

\begin{grammarlens}[title={What you've built}]
One more word for this chapter.
\end{grammarlens}

\subsection*{Your turn}
\begin{itemize}
\raggedright
  \item \emph{Say it:} \emph{wājib}
  \item \emph{Say it:} \emph{wājib}, once more
  \item \emph{Say it:} say \emph{wājib}
  \item \emph{Recall:} say the Arabic for a contract, then the Arabic for an exam, then say \emph{wājib} again
\end{itemize}

\begin{tcolorbox}[breakable,colback=teal!4,colframe=teal!35!black,title={Before you move on}]
What does \ar{واجب} mean? (\textquotedblleft{}Homework; a duty\textquotedblright{}.) Say it once more.
\end{tcolorbox}
